\\documentclass[11pt,a4paper]{article}
\\usepackage{fontspec}
\\usepackage{amsmath,amssymb,mathtools}
\\usepackage{geometry}
\\usepackage{microtype}
\\usepackage{unicode-math}
\\usepackage{polyglossia}
\\setmainlanguage{french}
\\setmainfont{Latin Modern Roman}
\\setmathfont{Latin Modern Math}
\\geometry{margin=2.2cm}

\\title{Aurore — Test de génération PDF}
\\author{}
\\date{}

\\begin{document}
\\maketitle

\\section{Test scientifique}

Soit $f : \\mathbb{R} \\to \\mathbb{R}$ définie par
\\[
f(x)=e^x.
\\]

On vérifie notamment une fraction :
\\[
\\frac{1}{x+1}+\\frac{2}{x-1}
=\\frac{3x+1}{x^2-1}.
\\]

Une limite :
\\[
\\lim_{x\\to +\\infty} e^x=+\\infty.
\\]

Une intégrale :
\\[
\\int_0^1 x^2\\,dx=\\frac{1}{3}.
\\]

Une matrice :
\\[
A=
\\begin{pmatrix}
1 & 2\\\\
3 & 4
\\end{pmatrix}.
\\]

Un ensemble :
\\[
E=\\left\\{x\\in\\mathbb{R}\\mid x^2<4\\right\\}.
\\]

\\section{Texte et Unicode}

Voici du texte français avec des accents : élève, équation, démonstration,
réussite, propriété, intégrale et géométrie.

\\section{Exercice}

Étudier les variations de $f(x)=e^x$ et justifier le résultat.

\\end{document}
